This paper addressed a subtle but high-impact failure mode in hexapod MPC: the optimizer could converge to physically inconsistent tarsus branches even while producing task-space-plausible trajectories. We showed that the problem was not primarily caused by gait timing, weight tuning, or runtime guard logic, but by insufficient solver-level enforcement of branch admissibility in the original SQP pipeline. To resolve this issue, we introduced a hard \emph{TarsusBranchConstraint} and solved the resulting optimal control problem with an Interior-Point Method.

The resulting controller eliminates branch-guard interventions in the reported experiments, restores stable forward progression, improves tracking accuracy, and remains comfortably within real-time computational limits. In addition to the short-horizon A/B comparison against the SQP baseline, the controller also sustained stable medium- and long-horizon endurance walking while preserving bounded yaw drift, limited base-height variation, and sub-5-ms average solve time. Beyond the immediate performance gain, the more important outcome is architectural: physical branch validity is now enforced inside the optimization problem rather than repaired after optimization. This shifts the controller from reactive damage control to proactive feasible-set design, which is essential for dependable legged locomotion.

The broader implication is that high-dimensional multi-legged systems require solver formulations that preserve mechanism semantics, not merely objective convergence. In such systems, deceptive local minima can appear numerically benign while remaining physically invalid. Our results suggest that IPM combined with explicit mechanism-aware hard constraints offers a promising direction for avoiding these pitfalls and for building MPC controllers with stronger safety interpretations.

Taken together, the paper shows that solver-visible branch-consistency constraints are not merely a numerical convenience, but a prerequisite for reliable online hexapod MPC. By pairing these constraints with an IPM backend and by documenting the engineering conditions required for reproducible real-time performance, we provide a controller design that is both scientifically interpretable and practically deployable.
